% ---------------------------------------------------------------------
%  07-operating.tex -- Running HPB: the operating guide
%  Practices, intake and typing, cadence, roles, coexistence with
%  existing frameworks, adoption path, anti-patterns.
% ---------------------------------------------------------------------
\section{Running HPB: The Operating Guide}
\label{sec:operating}

The model says what a correct HPB graph is, and the notation says how to
draw it. Neither says who changes the graph, when, and on what evidence.
This section does. The operating guide is deliberately small: a handful
of practices, one intake convention, four recurring meetings, roles that
already exist, and an adoption path. It is meant to be started by one
product manager on one product line and to scale by repetition, not by
adding ceremony. It prescribes \emph{what must be recorded} at each
decision; how engineering organizes its own work stays with whatever
team process is already in place (\cref{subsec:method-coexistence}).

\subsection{Four principles}
\label{subsec:method-principles}

\begin{enumerate}[label=\textbf{M\arabic*},leftmargin=*,itemsep=2pt]
  \item \textbf{Record at the decision, not after it.} Every change to the
    graph matches a decision that is already being taken: accepting a
    request, committing it to a release, publishing a version, dropping
    work. HPB attaches the record to that moment, so the graph is a
    by-product of deciding, not an extra reporting duty.
  \item \textbf{One accountable owner per node.} Every line, baseline and
    branch has exactly one accountable role. Many people contribute; one
    answers for it.
  \item \textbf{Closure is explicit and evidenced.} Nothing disappears.
    A branch either merges into a named baseline or is discarded with a
    written reason (\cref{rule:closure}). This is the principle that
    produces traceability, and the one list-based trackers break most
    easily, because there an item is simply set to ``closed''.
  \item \textbf{History is fixed, the plan is not.} Everything left of NOW
    is a record and is corrected only through audited events. Everything
    right of NOW is a plan and is expected to change. Plan changes should
    be cheap; history changes should be visible.
\end{enumerate}

\subsection{Practices}
\label{subsec:method-practices}

\Cref{tab:practices} lists the practices. Each one has a trigger, the
roles involved, the rule it enforces and what it leaves in the graph.
The point of the table is that the rules are not checked in a separate
compliance step: each rule is enforced at the one moment in the workflow
where it can be broken.

\begin{table}[tbp]
  \centering
  \small
  \caption{HPB practices. Each practice runs at a clear trigger and enforces
    at least one rule (R1--R7). Roles: A accountable, R responsible (if
    different), C consulted or approving. PM product manager; RM release or
    program manager; CCB change control board; Eng R\&D/engineering; QA
    quality assurance or laboratory; QR quality or regulatory function.}
  \label{tab:practices}
  \begin{tabularx}{\linewidth}{@{}>{\raggedright\arraybackslash}p{2.0cm}
                                  >{\raggedright\arraybackslash}X
                                  >{\raggedright\arraybackslash}p{2.2cm}
                                  >{\raggedright\arraybackslash}p{1.0cm}
                                  >{\raggedright\arraybackslash}X@{}}
    \toprule
    \textbf{Practice} & \textbf{Trigger} & \textbf{Roles} & \textbf{Rules} & \textbf{Leaves in the graph} \\
    \midrule
    Fork at intake
      & A request, defect report, task, lab evaluation or supplier notice is accepted at triage.
      & A: PM; R: support, Eng, QA
      & R1, R6
      & Typed branch with origin baseline (and configuration), opening date, owner. \\
    \addlinespace
    Target at planning
      & Release planning for a planned baseline; a priority change.
      & A: PM; C: RM
      & R4, R7
      & Target baseline; dashed continuation; earlier targets kept as history. \\
    \addlinespace
    Verify
      & A candidate of the target (build, prototype, pilot lot) is tested.
      & A: QA; R: Eng
      & R4, R7
      & Pass/fail step citing the candidate; a pass makes the branch verified. \\
    \addlinespace
    Gate review
      & Freeze or target date of a planned baseline or launch bundle.
      & A: RM; C: PM, Eng, QA
      & R4
      & Decision on the baseline: release, slip with new date, or move/discard the unverified members. \\
    \addlinespace
    Merge at release
      & The baseline is published (build released, or change order approved with production effectivity).
      & A: RM; R: Eng
      & R2, R3
      & Merge of every verified member; release notes derived. \\
    \addlinespace
    Discard with reason
      & Not viable, duplicate, not reproducible, superseded, out of scope; or resolved without a product baseline.
      & A: PM; R: owner; C: QR if safety or regulatory
      & R2, R7
      & Discard terminal with reason class, category and text. \\
    \addlinespace
    Reopen (audited)
      & A closure was recorded in error.
      & A: PM; C: second approver
      & R2, R7
      & Reopen event with reason and actors; original terminal kept. \\
    \addlinespace
    Line fork
      & A hardware change fails the interchangeability test.
      & A: CCB; R: PM, hardware Eng
      & R5
      & Diamond baseline, child line with parent link, change-order reference, parent life-cycle milestones. \\
    \addlinespace
    Record step
      & A sub-task completes, a test runs, an effectiveness check is made.
      & A: Eng or QA
      & R6, R7
      & Step (done/pending, pass/fail) with evidence attached. \\
    \bottomrule
  \end{tabularx}
\end{table}

\paragraph{Fork at intake.}
A request becomes a branch only when it is accepted at triage. At that
moment three facts are fixed: its type, its origin baseline and its
opening date (\cref{rule:origin}). For an issue, the origin is the
version in which the defect was \emph{seen}: the firmware version read
from the device, the software build, the hardware revision on the label.
It is mandatory. If nobody knows it, clarify before accepting, or accept
it with the origin flagged as \emph{inferred}; when the defect depends on
a combination of hardware and firmware, record the configuration too.
Triage is also where a request can be rejected before it ever becomes a
branch, which keeps the graph free of noise.

\paragraph{Granularity and typing.}
Diagrams are only comparable between product managers if the same kinds
of request become branches with the same types. \Cref{tab:intake} is the
default convention. A branch is a \emph{release-note-worthy unit of
change}: an epic or feature, a customer-visible defect, a lab campaign.
User stories, sub-tickets and commits stay on team boards; they may
appear on a branch as steps (\cref{rule:typed}). Security, regulatory
relevance, safety or a specific customer are \emph{attributes} of a
branch, not extra types, so the four lanes stay stable.

\begin{table}[tbp]
  \centering
  \small
  \caption{Default intake and typing convention. A request becomes one
    branch of the listed type; attributes are set in addition.
    Organizations may refine the table, but should publish it with their
    diagrams.}
  \label{tab:intake}
  \begin{tabularx}{\linewidth}{@{}>{\raggedright\arraybackslash}X
                                  >{\centering\arraybackslash}p{1.0cm}
                                  >{\raggedright\arraybackslash}p{5.6cm}@{}}
    \toprule
    \textbf{Request} & \textbf{Type} & \textbf{Origin baseline; attributes} \\
    \midrule
    New or changed customer-visible capability (epic, feature) & $\mathsf{I}$ & latest release; customer if requested by one \\
    Defect seen in a released version & $\mathsf{X}$ & version where it was seen (mandatory); configuration if relevant \\
    Security vulnerability & $\mathsf{X}$ & version where it was seen; attribute \emph{security} \\
    Regulatory or certification change & $\mathsf{T}$ & latest release; attribute \emph{regulatory}; $\mathsf{I}$ if it adds a capability \\
    Customer-specific request & $\mathsf{I}$ or $\mathsf{X}$ & as above; attribute \emph{customer}, with the commitment made \\
    Cost-down redesign, component substitution & $\mathsf{L}$ & latest release; line fork only if not interchangeable \\
    Supplier change or end-of-life notice & $\mathsf{L}$ or $\mathsf{T}$ & latest release; $\mathsf{L}$ to qualify an alternative, $\mathsf{T}$ for a last-time buy \\
    Field return with failure analysis & $\mathsf{X}$ & version of the returned unit; RMA evidence as steps \\
    Internal work (certificates, tooling, documentation) & $\mathsf{T}$ & latest release \\
    User story, sub-ticket, single test run & -- & a step of its branch, not a branch \\
    \bottomrule
  \end{tabularx}
\end{table}

\paragraph{Target at planning.}
Targeting commits an open branch to a planned version. It is the only
practice that is expected to repeat for the same branch: when priorities
change, the branch is re-targeted and the old target stays visible as a
faded dashed line (\cref{fig:discard}). My rule of thumb: target a branch
only when its owner can state what ``done'' means for the merge. A target
without an acceptance condition is a wish, and wishes inflate the gate.
Before committing the targeted set, compare its size with what recent
releases of the same kind actually contained (\cref{tab:metrics}), and check
that every dependency of a member is targeted at the same or an earlier
version.

\paragraph{Gate review.}
At the freeze date (or the target date, for teams without a freeze) the
release manager, with product, engineering and quality, goes through the
targeted set. Each member is verified, not yet verified, or no longer
wanted; each dependency is satisfied or not. \Cref{rule:gate} forbids
publishing while any member is unverified or depends on unreleased work,
so the review must end in one of three recorded outcomes:
\emph{release}, if everything is verified; \emph{slip}, with a new date
and the reason written on the baseline; or \emph{descope}, moving or
discarding the unverified members. The rule does not choose for you. It
forces the choice to be made and written down. \Cref{fig:release-gate}
shows the situation the review resolves.

\paragraph{Discard with reason.}
Discarding is a normal outcome, not a failure. The engineering-change
literature treats the rejection of proposed changes as part of the change
process \citep{jarratt2011ec}; HPB makes those decisions visible instead of
hiding them in a generic ``closed''. The reason is picked from the two
closed lists of \cref{sec:model-closing}, which makes discards countable,
plus free text and, for a duplicate, a link to the surviving branch, which
makes them auditable. Discarding a branch marked \emph{safety} or
\emph{regulatory} also needs the approval of the quality or regulatory
function.

\paragraph{The other practices.}
\emph{Verify}: whoever ran the test records the passed step with the
candidate it ran on (a release-candidate build, a prototype or a pilot
lot); that record replaces the familiar ``yes, it's in the release''.
\emph{Merge at release}: publishing a version merges all its verified
targeted branches at once, and a hardware baseline cites its change
order. \emph{Reopen} needs a second approver and keeps the original
terminal; a fix that regresses later is a \emph{new} issue from the
version where the regression is seen. \emph{Line fork}: when the change
control board approves a new reference, the parent gets dated life-cycle
milestones at the same time; leaving them open is the surest way to
accumulate maintenance nobody planned for. \emph{Record step}: sub-tasks,
test runs and effectiveness checks are steps, so a lab branch becomes an
evidence trail (\cref{fig:hardware-fork,fig:rma}).

\subsection{Cadence}
\label{subsec:method-cadence}

The practices are event-driven, but most events land in four recurring
meetings. Adapt the frequencies to the release rhythm of your product;
keep the order.

\begin{description}[leftmargin=0pt,itemsep=3pt,font=\normalfont\bfseries]
  \item[Weekly triage (30--45 minutes).] New requests from customers,
    support, sales, engineering and the lab. Practices: fork at intake;
    discard obvious duplicates. Read the diagram below the line first (new
    issues and their origins often reveal a bad release early), then above
    it. End by reviewing open, untargeted branches older than an agreed
    threshold (age of open work, \cref{tab:metrics}).
  \item[Release planning (per planned version).] Open branches, capacity
    from the team process, dependencies, customer commitments. Practice:
    target at planning, with the size and dependency checks. Output: the
    targeted set of the next version or launch bundle and, optionally,
    provisional targets for the one after. Planning only targets work that
    already exists as a branch, so every planned change has a recorded
    origin.
  \item[Gate review (at freeze or target date).] Gate review, then merge
    at release if the decision is to ship. The diagram keeps it short: the
    dashed lines entering the gate without a passed step are, by
    construction, exactly the open questions.
  \item[Line review (monthly or after each release).] Release
    contents, escaped defects, installed base and rollout where available,
    unported fixes, and the metrics of \cref{sec:metrics}. Decisions:
    life-cycle milestones of maintained lines and streams, whether a patch
    is needed, which fixes must be ported, and whether recent discard
    reasons point to a problem in intake quality.
\end{description}

\subsection{Roles}
\label{subsec:method-roles}

HPB introduces no new job titles. It gives accountability for graph
elements to roles that product organizations already have
(\cref{tab:practices}).

\begin{itemize}[leftmargin=*,itemsep=2pt]
  \item \textbf{Product manager.} Owns the product lines and the content of
    planned versions: triage, typing, targeting, discards. Can delegate,
    as long as the actor of every transition is recorded.
  \item \textbf{Release or program manager.} Owns dates and gates: runs the
    gate review and publishes versions and launch bundles.
  \item \textbf{Change control board.} Where one exists, approves change
    orders, and therefore hardware baselines and line forks; the HPB
    baseline references the approved order.
  \item \textbf{R\&D / engineering.} Owns improvement and task branches
    once accepted, and records their steps.
  \item \textbf{Quality / laboratory.} Owns lab branches, verification
    steps and their evidence; is consulted at every gate. The quality or
    regulatory function approves discards of safety- or
    regulation-relevant issues and the closure of corrective actions.
  \item \textbf{Support and field service.} The main source of issues.
    Owns the quality of the reported version and configuration, which is
    what escaped-defect counts depend on, and the installed-base data
    where it is kept.
\end{itemize}

\noindent In a small organization one person wears several hats. That is
fine, as long as the role under which each action is taken is visible in
the audit trail. Two segregations of duty stay even then: reopens and
safety- or regulation-relevant discards need a second approver.

\subsection{Living with Scrum, Kanban, Stage-Gate and PLM}
\label{subsec:method-coexistence}

HPB is a \emph{product-level} view. It does not schedule engineers, set
team WIP limits or replace the boards teams use every day. It sits above
them and answers what they are not designed to answer: from which product
state work started, which version closed it, and what else was in flight
on the same product at the same time. \Cref{tab:coexistence} gives the
mapping; \cref{fig:same-data} shows the same items in both views.

\begin{table}[tbp]
  \centering
  \small
  \caption{How HPB maps onto the units of common frameworks. HPB adds the
    product-state anchors (origin and closing baseline) without replacing
    the framework's own planning unit.}
  \label{tab:coexistence}
  \begin{tabularx}{\linewidth}{@{}>{\raggedright\arraybackslash}p{2.0cm}
                                  >{\raggedright\arraybackslash}X
                                  >{\raggedright\arraybackslash}X@{}}
    \toprule
    \textbf{Framework} & \textbf{Its unit} & \textbf{In HPB} \\
    \midrule
    Scrum
      & Product backlog item; sprint; increment \citep{schwaber2020scrum}.
      & A backlog item is (part of) a branch. A sprint is a team time-box,
        \emph{not} a baseline; an increment becomes a baseline only when it
        is released to customers. \\
    \addlinespace
    Kanban
      & Card; column; WIP limit \citep{anderson2010kanban}.
      & A card is a branch or a step of one. Columns refine the HPB states
        open and targeted; ``done'' becomes a merge only once the closing
        version is published. \\
    \addlinespace
    Stage-Gate
      & Project; stage; go/kill gate \citep{cooper1990stagegate}.
      & A new-reference project is lab and improvement work ending in a
        hardware baseline and a line fork. A \emph{kill} is a discard with
        reason. Stage-Gate gates decide whether to keep investing; HPB
        gates decide whether a release is complete. \\
    \addlinespace
    PLM / CM
      & Change order; configuration baseline \citep{stark2015plm,iso10007}.
      & A released HPB baseline references the change order or
        configuration baseline that formally defines it. \\
    \bottomrule
  \end{tabularx}
\end{table}

The most common confusion is the sprint: it is a team time-box, not a
baseline, because firmware and hardware ship less often than teams finish
sprints. In practice the product owner usually holds the HPB product
manager role, and board and diagram share the same tracker items, so
there is no double entry.

\paragraph{Regulated contexts.}
HPB does not replace configuration management as defined in IEEE~828 and
ISO~10007 \citep{ieee828,iso10007}, nor change processes in a PLM system
\citep{stark2015plm}; it adds the planning view and the link between the
intent of a change (the branch) and the configuration that realized it
(the baseline). Used as evidence, it needs at least append-only history,
identified actors and approvers on every closing decision, links to the
requirement, risk and test records, and a per-version export. Whether
that satisfies a particular regulation is outside the scope of this paper.

\subsection{Start on Monday: an adoption path}
\label{subsec:method-adoption}

HPB needs no migration and no new tool on day one. This is the sequence I
recommend.

\paragraph{Week 1: one line, drawn by hand.}
Pick \emph{one} product line with parallel work and a release coming up,
not the whole portfolio. Write down its last three to five released
versions and where the line came from. Go through the open items in your
trackers and record, for each release-note-worthy one, its type and its
best-known origin, marking guesses as \emph{inferred}. Create the next
planned version and target what is already promised for it: the gate is
now visible. Draw it, on a whiteboard, with \texttt{hpb.sty}
(\cref{app:package}) or in a few lines of HPB Markup (\cref{sec:tooling}),
and show it at your next triage.

\paragraph{Month 1: run the cadence.}
Apply \emph{fork at intake} to every new request: no branch without a type
and an origin. Run weekly triage on the diagram, issues first. Run the
first gate review at the next freeze and record its outcome. When the
version ships, merge its verified branches; from then on, release notes
are generated, not written. Publish your intake convention
(\cref{tab:intake}) with your changes.

\paragraph{Quarter 1: extend and measure.}
After a full release cycle, add the next lines, with their hardware
revisions, line forks and maintained streams. Set life-cycle milestones
for every old reference still maintained. Hold the first line review with
the metrics of \cref{sec:metrics}, and connect HPB to your tracker so the
diagram is generated from the items, not maintained by hand.

\subsection{Anti-patterns}
\label{subsec:method-antipatterns}

These are easy to spot in the diagram, and each one tells you which
practice is being skipped.

\begin{itemize}[leftmargin=*,itemsep=1pt]
  \item \emph{Silent closure.} A tracker item is closed but its branch has
    no merge or discard terminal (breaks \cref{rule:closure}).
  \item \emph{Eternal branch.} An open, untargeted branch that keeps
    outliving releases: usually a discard nobody dared to decide.
  \item \emph{Gate stuffing.} A targeted set much larger than recent
    releases, which turns every gate review into a slip.
  \item \emph{Floating evidence.} Test reports and decisions kept outside
    the graph (breaks \cref{rule:discussion}).
  \item \emph{In-place reference edit.} A non-interchangeable hardware
    change without a new line (breaks \cref{rule:lineage}).
  \item \emph{Unported fix.} A merged issue known to affect another
    maintained stream, with no port and no decision not to port.
  \item \emph{Shipping ahead of a dependency.} A dependency connector that
    points to the right of the gate (blocked by \cref{rule:gate}).
\end{itemize}
